% Seccion "Introduction" del paper de bin packing 3D (BSG_BPP).
\section{Introduction}
\label{sec:intro}

The three-dimensional bin packing problem (3D-BPP) asks for the minimum number of identical bins that
hold a set of rectangular boxes \citep{Martello2000}. It appears whenever goods are loaded into
containers, trucks or pallets, and in practice it comes with constraints such as the vertical support
of the boxes \citep{Bortfeldt2013}. The problem is strongly NP-hard, exact methods solve only small
instances, and the benchmark of \citet{Martello2000} has been the test bed of two decades of
heuristics.

Two families of heuristics dominate the benchmark. Metaheuristics such as guided local search
\citep{Faroe2003}, tabu search \citep{Crainic2009}, GRASP with variable neighbourhood descent
\citep{Parreno2010} and the biased random key genetic algorithm (BRKGA) of \citet{Goncalves2013} search
over complete solutions. Column generation methods \citep{Mahvash2018,Elhedhli2019} build bins as
patterns with a pricing problem guided by dual prices and choose the best combination of patterns with
a set covering or partitioning model. The first family produces many good bins during the search but
keeps only the solution in which they appear; the second recombines bins but generates each one with
a costly pricing problem.

This paper combines the cheap generation of the first family with the recombination of the second.
Each bin is loaded by BSG, a beam search for the single container loading problem
\citep{Araya2014,Araya2017}, guided by profits per box type. A squeaky wheel loop \citep{Joslin1999}
repeats the construction and raises the profits of the box types that ended in the least loaded bins.
Every bin that the search builds is stored in a pool of patterns, together with the best complete
loads that BSG explores and discards (the \emph{harvest}). After the search, a set covering model over
the pool chooses the combination of bins that loads all the boxes with the fewest bins; when removing a
box can make a bin infeasible, as with vertical support, a set partitioning model is used instead. The
LP relaxation of the model bounds what the pool can offer: in most runs it proves that the incumbent is
the best combination in the pool and no integer program is solved.

\paragraph{Contributions.}
\begin{itemize}
\item A generator of bin patterns that uses a single-bin loader inside a squeaky wheel loop over the
  profits of the box types, and that harvests, at almost no cost, the complete loads that the loader
  explores and discards (Sections~\ref{sec:m-construct}--\ref{sec:m-pool}).
\item A recombination of the pool by set covering, or by set partitioning under vertical support, with
  a stopping rule and a reduced-cost fixing derived from the LP relaxation of the pool
  (Section~\ref{sec:m-cover}).
\item An experimental study on the benchmark of \citet{Martello2000}, with fixed orientation and with
  rotations, and on industrial instances with vertical support \citep{Elhedhli2019}, with an ablation
  of the components, matched-time controls and instances held out from the design
  (Section~\ref{sec:experiments-bpp}).
\end{itemize}

\paragraph{Summary of the results.}
Without rotations, a configuration that runs 6.7~s per instance obtains totals at the level of BRKGA
\citep{Goncalves2013}, and a configuration of 100~s per instance obtains 7249.3 bins on classes 1 and
4--8, against 7252 for the column generation of \citet{Mahvash2018}. With rotations, the longer
configuration obtains 8982.3 bins on the eight classes, against 9013 for \citet{Mahvash2018}. The
recombination saves 10 to 17 bins with respect to the squeaky wheel loop with the same time on
instances that were not used to design the method, and about 95\% of the runs end with a proof that the pool holds no better solution,
in most of them by the LP alone. On industrial instances with 70\% support on every box,
the construction alone uses 6.0\% fewer bins than the column generation over layers of \citet{Elhedhli2019} when $n\ge 500$, although
the recombination did not improve any solution there.

\paragraph{Outline.}
Section~\ref{sec:related} reviews the related work and Section~\ref{sec:problem} defines the problem.
Section~\ref{sec:method} describes the method and Section~\ref{sec:experiments-bpp} the experiments.
Section~\ref{sec:conclusions} concludes.
